\documentclass[10pt,a4paper]{amsart}
\usepackage{amsmath,amssymb,amsthm,graphicx}
% Retained minimal preamble for the Stacks Project chapter sites-cohomology.tex (Cohomology on Sites).
% Verbatim excerpt of the chapter preamble: Xy-pic setup (native preamble lines 14-18) and the native
% theorem declarations with the mathematical macros (native preamble lines 154-217). Copied mechanically;
% nothing added, removed or redefined. The wrapper class and the amsmath/amssymb/amsthm/graphicx base
% packages are supplied by the assembler, so the class/externaldocument part of the native preamble is omitted.
% For commutative diagrams we use Xy-pic
\usepackage[all]{xy}

% We use 2cell for 2-commutative diagrams.
\xyoption{2cell}
% Theorem environments.
%
\theoremstyle{plain}
\newtheorem{theorem}[subsection]{Theorem}
\newtheorem{proposition}[subsection]{Proposition}
\newtheorem{lemma}[subsection]{Lemma}

\theoremstyle{definition}
\newtheorem{definition}[subsection]{Definition}
\newtheorem{example}[subsection]{Example}
\newtheorem{exercise}[subsection]{Exercise}
\newtheorem{situation}[subsection]{Situation}

\theoremstyle{remark}
\newtheorem{remark}[subsection]{Remark}
\newtheorem{remarks}[subsection]{Remarks}

\numberwithin{equation}{subsection}

% Macros
%
\def\lim{\mathop{\mathrm{lim}}\nolimits}
\def\colim{\mathop{\mathrm{colim}}\nolimits}
\def\Spec{\mathop{\mathrm{Spec}}}
\def\Hom{\mathop{\mathrm{Hom}}\nolimits}
\def\Ext{\mathop{\mathrm{Ext}}\nolimits}
\def\SheafHom{\mathop{\mathcal{H}\!\mathit{om}}\nolimits}
\def\SheafExt{\mathop{\mathcal{E}\!\mathit{xt}}\nolimits}
\def\Sch{\mathit{Sch}}
\def\Mor{\mathop{\mathrm{Mor}}\nolimits}
\def\Ob{\mathop{\mathrm{Ob}}\nolimits}
\def\Sh{\mathop{\mathit{Sh}}\nolimits}
\def\NL{\mathop{N\!L}\nolimits}
\def\CH{\mathop{\mathrm{CH}}\nolimits}
\def\proetale{{pro\text{-}\acute{e}tale}}
\def\etale{{\acute{e}tale}}
\def\QCoh{\mathit{QCoh}}
\def\Ker{\mathop{\mathrm{Ker}}}
\def\Im{\mathop{\mathrm{Im}}}
\def\Coker{\mathop{\mathrm{Coker}}}
\def\Coim{\mathop{\mathrm{Coim}}}

% Boxtimes
%
\DeclareMathSymbol{\boxtimes}{\mathbin}{AMSa}{"02}

%
% Macros for moduli stacks/spaces
%
\def\QCohstack{\mathcal{QC}\!\mathit{oh}}
\def\Cohstack{\mathcal{C}\!\mathit{oh}}
\def\Spacesstack{\mathcal{S}\!\mathit{paces}}
\def\Quotfunctor{\mathrm{Quot}}
\def\Hilbfunctor{\mathrm{Hilb}}
\def\Curvesstack{\mathcal{C}\!\mathit{urves}}
\def\Polarizedstack{\mathcal{P}\!\mathit{olarized}}
\def\Complexesstack{\mathcal{C}\!\mathit{omplexes}}
% \Pic is the operator that assigns to X its picard group, usage \Pic(X)
% \Picardstack_{X/B} denotes the Picard stack of X over B
% \Picardfunctor_{X/B} denotes the Picard functor of X over B
\def\Pic{\mathop{\mathrm{Pic}}\nolimits}
\def\Picardstack{\mathcal{P}\!\mathit{ic}}
\def\Picardfunctor{\mathrm{Pic}}
\def\Deformationcategory{\mathcal{D}\!\mathit{ef}}
\usepackage{cleveref}
\crefname{originallemma}{original lemma}{original lemmas}
\Crefname{originallemma}{Original lemma}{Original lemmas}
\makeatletter
\newlabel{sites-definition-site}{{00VH}{source}{}{sourcepointer-1}{}}
\newlabel{sites-definition-site@cref}{{[originallemma][0][]00VH}{[1][0][]source}{}{}{}}
\makeatother
\crefname{originallemma}{original lemma}{original lemmas}
\Crefname{originallemma}{Original lemma}{Original lemmas}
\makeatletter
\newlabel{sites-lemma-directed-colimits-sections}{{0738}{source}{}{sourcepointer-2}{}}
\newlabel{sites-lemma-directed-colimits-sections@cref}{{[originallemma][0][]0738}{[1][0][]source}{}{}{}}
\makeatother
\crefname{originallemma}{original lemma}{original lemmas}
\Crefname{originallemma}{Original lemma}{Original lemmas}
\makeatletter
\newlabel{injectives-theorem-sheaves-injectives}{{01DP}{source}{}{sourcepointer-3}{}}
\newlabel{injectives-theorem-sheaves-injectives@cref}{{[originallemma][0][]01DP}{[1][0][]source}{}{}{}}
\makeatother
\crefname{originallemma}{original lemma}{original lemmas}
\Crefname{originallemma}{Original lemma}{Original lemmas}
\makeatletter
\newlabel{algebra-lemma-directed-colimit-exact}{{00DB}{source}{}{sourcepointer-4}{}}
\newlabel{algebra-lemma-directed-colimit-exact@cref}{{[originallemma][0][]00DB}{[1][0][]source}{}{}{}}
\makeatother
\crefname{originallemma}{original lemma}{original lemmas}
\Crefname{originallemma}{Original lemma}{Original lemmas}
\makeatletter
\newlabel{lemma-cech-vanish-collection}{{03F9}{source}{}{sourcepointer-5}{}}
\newlabel{lemma-cech-vanish-collection@cref}{{[originallemma][0][]03F9}{[1][0][]source}{}{}{}}
\makeatother
\crefname{originallemma}{original lemma}{original lemmas}
\Crefname{originallemma}{Original lemma}{Original lemmas}
\makeatletter
\newlabel{lemma-injective-trivial-cech}{{03AW}{source}{}{sourcepointer-6}{}}
\newlabel{lemma-injective-trivial-cech@cref}{{[originallemma][0][]03AW}{[1][0][]source}{}{}{}}
\makeatother
\title{Filtered colimits commute with cohomology for a collection of finite coverings cofinal in each object}
\author{The Stacks Project Authors}
\date{stacks-project:sites-cohomology.tex:local-snapshot-20261011; GNU FDL 1.2 or later}
\begin{document}
\maketitle
\noindent\textit{Editorial note.} All delivered mathematical text is the author's own, copied from the pinned Stacks Project chapter file sites-cohomology.tex; no statement or proof step was rewritten, reordered or removed.\par\medskip
\noindent\textit{Editorial note.} The proof is the author's complete argument: the reduction to the vanishing of the cohomology of a filtered colimit of injectives, and the application of the cofinality criterion, are both carried through inside the counted span.\par\medskip
\noindent\textit{Editorial note.} Every labelled reference occurring inside the counted spans points to a chapter of the Project that is not part of this card's source tree. Each is carried in assembly.external\_references as a non-load-bearing navigation pointer whose source location is the invocation inside the counted span, so no load-bearing mathematics is replaced by a stub.\par\medskip
\noindent\textit{Editorial note.} The checked chapter list for same-name collisions is the set of stacks chapters staged in this repository: chow, cohomology, divisors, etale-cohomology, flat, more-algebra, more-morphisms, morphisms, spaces-more-morphisms, varieties together with this card's chapter sites-cohomology.\par\medskip
\noindent\textit{Editorial note.} The preamble is the chapter's own retained minimal preamble (sources/stacks-sites-cohomology/retained-preamble.tex), a verbatim excerpt of the native preamble: the Xy-pic setup and the native theorem declarations with the mathematical macros.\par\medskip
\noindent\textit{Allowed standard prerequisites (editorial).}\begin{itemize}
\item Cech cohomology of a finite covering of a site. Same chapter, the Cech complex of a covering and the comparison between Cech cohomology and cohomology for a cofinal system of coverings. Supplies the coverings used to compute the cohomology of the colimit and the criterion applied at the end of the proof.
\item Functorial injective embeddings in Ab(C). Injectives, the existence of functorial injective resolutions in a Grothendieck abelian category, applied to abelian sheaves on a site. Supplies the construction of the filtered system of injective sheaves I\_i and their colimit I used to reduce the statement to the vanishing of H\textasciicircum{}p on a filtered colimit of injectives.
\item Vanishing criterion for a filtered colimit of injective sheaves over the collection Cov. Same chapter, the lemma stating that the cohomology of a filtered colimit of injective sheaves vanishes in degrees at least one when the collection of coverings is cofinal. Supplies the final step: the claim that the degree-p part vanishes is proved by invoking that criterion with the data at hand.
\end{itemize}\medskip
\ifdefined\hypertarget\hypertarget{sourcepointer-1}{}\fi
\noindent\textit{Source reference.} Original lemma 00VH: sources/stacks-sites-cohomology/pinned-original-source.tex, source lines 2521--2521. Navigation pointer invoked at proof line 2521; the counted span contains the pointer only and the cited result is not reproduced inside it, so the pointer is registered as a non-load-bearing navigation stub in the established form for Stacks material. No bibliography citation is involved.\par\medskip
\ifdefined\hypertarget\hypertarget{sourcepointer-2}{}\fi
\noindent\textit{Source reference.} Original lemma 0738: sources/stacks-sites-cohomology/pinned-original-source.tex, source lines 2514--2514. Navigation pointer invoked at proof line 2514; the counted span contains the pointer only and the cited result is not reproduced inside it, so the pointer is registered as a non-load-bearing navigation stub in the established form for Stacks material. No bibliography citation is involved.\par\medskip
\ifdefined\hypertarget\hypertarget{sourcepointer-3}{}\fi
\noindent\textit{Source reference.} Original lemma 01DP: sources/stacks-sites-cohomology/pinned-original-source.tex, source lines 2557--2557. Navigation pointer invoked at proof line 2557; the counted span contains the pointer only and the cited result is not reproduced inside it, so the pointer is registered as a non-load-bearing navigation stub in the established form for Stacks material. No bibliography citation is involved.\par\medskip
\ifdefined\hypertarget\hypertarget{sourcepointer-4}{}\fi
\noindent\textit{Source reference.} Original lemma 00DB: sources/stacks-sites-cohomology/pinned-original-source.tex, source lines 2572--2572. Navigation pointer invoked at proof line 2572; the counted span contains the pointer only and the cited result is not reproduced inside it, so the pointer is registered as a non-load-bearing navigation stub in the established form for Stacks material. No bibliography citation is involved.\par\medskip
\ifdefined\hypertarget\hypertarget{sourcepointer-5}{}\fi
\noindent\textit{Source reference.} Original lemma 03F9: sources/stacks-sites-cohomology/pinned-original-source.tex, source lines 1565--1565. Navigation pointer invoked at proof line 2606; the counted span contains the pointer only and the cited result is not reproduced inside it, so the pointer is registered as a non-load-bearing navigation stub in the established form for Stacks material. No bibliography citation is involved.\par\medskip
\ifdefined\hypertarget\hypertarget{sourcepointer-6}{}\fi
\noindent\textit{Source reference.} Original lemma 03AW: sources/stacks-sites-cohomology/pinned-original-source.tex, source lines 1334--1334. Navigation pointer invoked at proof line 2616; the counted span contains the pointer only and the cited result is not reproduced inside it, so the pointer is registered as a non-load-bearing navigation stub in the established form for Stacks material. No bibliography citation is involved.\par\medskip
\begin{lemma}
\label{lemma-colim-works-over-collection}
Let $\mathcal{C}$ be a site. Let $\text{Cov}_\mathcal{C}$ be the set
of coverings of $\mathcal{C}$ (see
Sites, Definition \ref{sites-definition-site}). Let
$\mathcal{B} \subset \Ob(\mathcal{C})$, and
$\text{Cov} \subset \text{Cov}_\mathcal{C}$
be subsets. Assume that
\begin{enumerate}
\item For every $\mathcal{U} \in \text{Cov}$ we have
$\mathcal{U} = \{U_i \to U\}_{i \in I}$ with $I$ finite,
$U, U_i \in \mathcal{B}$ and every
$U_{i_0} \times_U \ldots \times_U U_{i_p} \in \mathcal{B}$.
\item For every $U \in \mathcal{B}$ the coverings of $U$
occurring in $\text{Cov}$ is a cofinal system of coverings of $U$.
\end{enumerate}
Then the map
$$
\colim_i H^p(U, \mathcal{F}_i)
\longrightarrow
H^p(U, \colim_i \mathcal{F}_i)
$$
is an isomorphism for every $p \geq 0$, every $U \in \mathcal{B}$, and
every filtered diagram $\mathcal{I} \to \textit{Ab}(\mathcal{C})$.
\end{lemma}

\begin{proof}
To prove the lemma we will argue by induction on $p$.
Note that we require in (1) the coverings $\mathcal{U} \in \text{Cov}$
to be finite, so that all the elements of $\mathcal{B}$ are quasi-compact.
Hence (2) and (1) imply that any $U \in \mathcal{B}$ satisfies the hypothesis
of Sites, Lemma \ref{sites-lemma-directed-colimits-sections} (4).
Thus we see that the result holds for $p = 0$.
Now we assume the lemma holds for $p$ and prove it for $p + 1$.

\medskip\noindent
Choose a filtered diagram
$\mathcal{F} : \mathcal{I} \to \textit{Ab}(\mathcal{C})$,
$i \mapsto \mathcal{F}_i$.
Since $\textit{Ab}(\mathcal{C})$ has functorial injective embeddings, see
Injectives, Theorem \ref{injectives-theorem-sheaves-injectives},
we can find a morphism of filtered diagrams
$\mathcal{F} \to \mathcal{I}$
such that each $\mathcal{F}_i \to \mathcal{I}_i$ is an injective map of
abelian sheaves into an injective abelian sheaf. Denote $\mathcal{Q}_i$
the cokernel so that we have short exact sequences
$$
0 \to
\mathcal{F}_i \to
\mathcal{I}_i \to
\mathcal{Q}_i \to 0.
$$
Since colimits of sheaves are the sheafification of colimits on the level
of presheaves, since sheafification is exact, and since filtered
colimits of abelian groups are exact
(see Algebra, Lemma \ref{algebra-lemma-directed-colimit-exact}),
we see the sequence
$$
0 \to
\colim_i \mathcal{F}_i \to
\colim_i \mathcal{I}_i \to
\colim_i \mathcal{Q}_i \to 0.
$$
is also a short exact sequence. We claim that
$H^q(U, \colim_i \mathcal{I}_i) = 0$ for all $U \in \mathcal{B}$
and all $q \geq 1$. Accepting this claim
for the moment consider the diagram
$$
\xymatrix{
\colim_i H^p(U, \mathcal{I}_i) \ar[d] \ar[r] &
\colim_i H^p(U, \mathcal{Q}_i) \ar[d] \ar[r] &
\colim_i H^{p + 1}(U, \mathcal{F}_i) \ar[d] \ar[r] &
0 \ar[d] \\
H^p(U, \colim_i \mathcal{I}_i) \ar[r] &
H^p(U, \colim_i \mathcal{Q}_i) \ar[r] &
H^{p + 1}(U, \colim_i \mathcal{F}_i) \ar[r] &
0
}
$$
The zero at the lower right corner comes from the claim and the
zero at the upper right corner comes from the fact that the sheaves
$\mathcal{I}_i$ are injective.
The top row is exact by an application of
Algebra, Lemma \ref{algebra-lemma-directed-colimit-exact}.
Hence by the snake lemma we deduce the
result for $p + 1$.

\medskip\noindent
It remains to show that the claim is true. We will use
Lemma \ref{lemma-cech-vanish-collection}.
By the result for $p = 0$ we see that for $\mathcal{U} \in \text{Cov}$
we have
$$
\check{\mathcal{C}}^\bullet(\mathcal{U}, \colim_i \mathcal{I}_i)
=
\colim_i \check{\mathcal{C}}^\bullet(\mathcal{U}, \mathcal{I}_i)
$$
because all the $U_{j_0} \times_U \ldots \times_U U_{j_p}$
are in $\mathcal{B}$. By
Lemma \ref{lemma-injective-trivial-cech}
each of the complexes in the colimit of {\v C}ech complexes is
acyclic in degree $\geq 1$. Hence by
Algebra, Lemma \ref{algebra-lemma-directed-colimit-exact}
we see that also the {\v C}ech complex
$\check{\mathcal{C}}^\bullet(\mathcal{U}, \colim_i \mathcal{I}_i)$
is acyclic in degrees $\geq 1$. In other words we see that
$\check{H}^p(\mathcal{U}, \colim_i \mathcal{I}_i) = 0$
for all $p \geq 1$. Thus the assumptions of
Lemma \ref{lemma-cech-vanish-collection}.
are satisfied and the claim follows.
\end{proof}

\end{document}
